{\color{gray}\hrule}
\begin{center}
\section{Appendix}
\bigskip
\end{center}
{\color{gray}\hrule}
\vspace{0.5cm}
\begin{multicols}{2}

\subsection{Training Setup and Hyperparameters}

The training parameters were kept consistent across both the baseline replication and the novel anti-biomimetic regimen to ensure comparability[cite: 320, 938].

\begin{itemize}
\item \textbf{Optimizer:} Stochastic Gradient Descent (SGD) with a Nesterov momentum of 0.9[cite: 280, 938].

\item \textbf{Loss Function:} Categorical cross-entropy[cite: 280, 937].

\item \textbf{Learning Rate:} Constant 0.001[cite: 280, 935].

\item \textbf{Batch Size:} 128[cite: 280, 936].

\item \textbf{Epochs (Baseline):} 200 epochs, with the transition occurring at epoch 100[cite: 281, 921].

\item \textbf{Epochs (This Work):} 120 epochs, with the transition occurring at epoch 60[cite: 504, 932].

\end{itemize}

\subsection{Data Preprocessing}

Images from the ImageNet database were subjected to the following preprocessing pipeline before being fed into the network[cite: 282, 940]:

\begin{enumerate}
\item \textbf{Resizing and Cropping:} Original $256 \times 256$ pixel images were randomly cropped to $227 \times 227$ pixel segments for the AlexNet architecture[cite: 282, 942].

\item \textbf{Augmentation:} Random horizontal flipping was applied during training[cite: 282, 943].

\item \textbf{Normalization:} Pixel values were rescaled from a $[0, 255]$ distribution to a $[-1, 1]$ distribution[cite: 282, 944].

\item \textbf{Degradation (Blur):} When applying blur, a Gaussian kernel with a $\sigma = 4$ was utilized[cite: 284, 329].

\end{enumerate}

\subsection{Receptive Field Metrics}

To characterize the first-layer filters of the modified AlexNet, the following metrics were employed:

\subsubsection{Color Metric}

The colorfulness $m$ of a filter was determined by computing pixel-wise intensity differences across the RGB channels[cite: 286, 1002]:

\begin{equation}
x = R \cos(0^\circ) + G \cos(120^\circ) + B \cos(-120^\circ) [cite: 287, 1003]
\end{equation}

\begin{equation}
y = R \sin(0^\circ) + G \sin(120^\circ) + B \sin(-120^\circ) [cite: 288, 1004]
\end{equation}

\begin{equation}
m = \sqrt{x^2 + y^2} [cite: 289, 1005]
\end{equation}

The final metric was the average of the top 48 (approximately top 10\%) most colorful pixels in the $22 \times 22$ filter[cite: 294].

\subsubsection{Spatial Frequency Metric}

The spatial frequency was summarized as a weighted average frequency[cite: 298, 1007]:

\begin{equation}
\text{Weighted Average Frequency} = \frac{\sum (amp(f) \cdot f)}{\sum amp(f)} [cite: 300, 1008]
\end{equation}

where $amp(f)$ represents the amplitude of frequency $f$ obtained via a 2D-FFT and radial averaging[cite: 297, 301].

\end{multicols}
